\begin{titlepage}
\thispagestyle{empty}
\begin{tikzpicture}[remember picture,overlay]
\fill[BookInk] (current page.north west) rectangle ([xshift=14mm]current page.south west);
\draw[BookAccent,line width=1.5pt] ([xshift=30mm,yshift=-42mm]current page.north west)--([xshift=-25mm,yshift=-42mm]current page.north east);
\end{tikzpicture}
\vspace*{25mm}
{\sffamily\color{BookAccent}\large 纯粹数学研讨会论文集\quad 第 33 卷}\par
\vspace{22mm}
{\Huge\sffamily\bfseries\color{BookInk}自守形式、表示\par\vspace{5mm}与 $L$ 函数}\par
\vspace{12mm}
{\Large\itshape Automorphic Forms, Representations,\par and $L$-functions}\par
\vspace{16mm}
{\Large A. Borel\quad W. Casselman\quad 编}\par
\vspace{6mm}
{\large 第一、二册中文合订本}\par
\vfill
{\color{BookGray}原版：American Mathematical Society，1979\par
1977 年俄勒冈州立大学夏季研究讲习班论文集}\par
\vspace{13mm}
\end{titlepage}
\chapter*{原版出版说明}
\addcontentsline{toc}{chapter}{原版出版说明}
\begin{center}
\textbf{美国数学会纯粹数学研讨会论文集}\par
1977 年 7 月 11 日至 8 月 5 日\par
俄勒冈州科瓦利斯，俄勒冈州立大学\par\medskip
A. Borel 与 W. Casselman 编\par\medskip
由美国数学会编制，部分经费由美国国家科学基金会\par
MCS 76-24539 号资助提供。
\end{center}

\paragraph*{美国国会图书馆在版编目资料（Library of Congress Cataloging-in-Publication Data）}
纯粹数学研讨会，俄勒冈州立大学，1977 年。\par
《自守形式、表示与 $L$ 函数》。\par
（纯粹数学研讨会论文集，第 33 卷。）\par
含参考文献与索引。\par
主题：1. 自守形式——会议。2. Lie 群——会议。3. 群的表示——会议。4. $L$ 函数——会议。\par
附加著录：I. Borel, Armand。II. Casselman, W., 1941--。III. 美国数学会。IV. 书名。V. 丛书。\par
分类号：QA331.S937\quad 1977\quad 512'.7\quad 78-21184。\par
ISBN 0-8218-1435-4（第一册）。\par
ISBN 0-8218-1437-0（第二册）。\par
AMS（MOS）主题分类（1970）：主要分类 10D20、22E50、12A70；次要分类 20G25、20G35、14G10、22E55、22E45。

\paragraph*{原版版权说明}
版权 \textcopyright\ 1979，美国数学会（American Mathematical Society）。美国印刷。除授予美国政府的权利外，美国数学会保留一切权利。

第一册所载复印与再版说明如下：本出版物的个人读者，以及代其行事的非营利图书馆，可以对材料作合理使用，例如为教学或研究复印一章。书评可以引用本出版物的简短段落，但须按惯例注明来源。对本出版物任何材料（包括摘要）的再出版、系统复印或多份复制，仅可在获得美国数学会许可的情况下进行。申请许可请致函美国数学会出版助理，P. O. Box 6248, Providence, Rhode Island 02940-6248，或发送电子邮件至 \texttt{reprint-permission@ams.org}。

第一册并载：本书所用纸张为无酸纸，符合旨在确保永久性与耐久性的准则。

第二册所载说明：未经出版者许可，不得以任何形式复制本书。

\chapter*{版本说明}
\addcontentsline{toc}{chapter}{版本说明}
\providecommand{\DraftNotice}{}\DraftNotice
本书按原有的四个主题部分，将 \textit{Automorphic Forms, Representations, and $L$-functions} 第一、二册合编为一卷。各篇文章保留原作者署名、内部章节编号、定理编号、公式编号及参考文献标签；全书章号与译本页码连续编排。正文引用的原书页码仍按册号区分。

本次校订调整中文语序、句式与段落衔接，并核查译本中的数学陈述与公式。发现的问题在相应位置以“译者注（本次校订）”说明，给出修正方式或补充理由；原有译注继续保留。

第二册 Ihara《同余关系与 Shimura 曲线》的原书第 295 页，依据用户随后补充的页面图像译出并重排，与所供 PDF 中的相邻页面接续。

技术术语在每篇首次出现时附通行英文，专名保留拉丁字母。法文文章亦译为中文。论述中的“已知”“猜想”“目前”等表述沿用原文的历史语境；本译本不据后来的研究结果改写作者论断。

排版使用 XeLaTeX。公式以数学环境重排，交换图与可重建图表优先使用 TikZ、TikZ-CD 和原生表格绘制；确实无法还原的图像保留原图来源。原文如含程序代码，采用带语法高亮的代码环境。参考文献题名保留原文以便检索。
\chapter*{前言（Foreword）}
\addcontentsline{toc}{chapter}{前言（Foreword）}
美国数学会（AMS）第二十五届夏季研究讲习班以自守形式（automorphic forms）、表示（representations）与 $L$ 函数（$L$-functions）为主题，于 1977 年 7 月 11 日至 8 月 5 日在俄勒冈州科瓦利斯的俄勒冈州立大学举行，由美国国家科学基金会拨款资助。组织委员会由 A. Borel、W. Casselman（共同主席）、P. Deligne、H. Jacquet、R. P. Langlands 和 J. Tate 组成。本卷所收文章包括讲习班的讲义笔记，多数经过修订，另有少数后来撰写的文章。

讲习班的主要目标之一，是讨论与约化群（reductive groups）上的自守形式或约化群的自守表示（automorphic representations）相联系的 $L$ 函数，相关的局部与整体问题（local and global problems），以及它们与代数数论、代数几何中的 $L$ 函数之间的关系，例如 Artin $L$ 函数和 Hasse--Weil zeta 函数。这个广泛的主题可以追溯到 E. Hecke、C. L. Siegel 等人的工作；近几年来，这一领域进展显著，并仍在迅速发展。其中一个重要推动因素，是在阿代尔群（adelic groups）的框架内系统运用无限维表示（infinite-dimensional representations）。这些进展汇集了多个领域的方法，其中一些需要较多的预备知识。因此，除关于近期工作、当前研究及未解决问题的讨论班与讲座之外，讲习班还安排了具有较强入门性质的讲座，甚至系列讲座，内容涵盖约化群及其表示、数论等背景材料，并对若干相对简单的情形作了详尽讨论。

本卷文章分为四个主要部分，这一划分也大致对应于各部分所需的预备知识。第一部分讨论约化群的结构以及局部域（local fields）上约化群的无限维表示。其中五篇提供基本的背景材料，其余各篇讨论近期进展。第二部分讨论自守形式与自守表示，侧重分析理论。前四篇讨论有关的基本事实与定义，以及两种表述之间的联系。另有两篇讨论 Eisenstein 级数（Eisenstein series）和迹公式（trace formula），先讨论 $\GL_2$，再讨论更一般的情形。事实上，迹公式和轨道积分（orbital integrals）是整个讲习班反复出现的主题，也见于其他部分的若干篇文章。最后四篇的主要主题，是将亚辛群（metaplectic group）的振子表示（oscillator representation）限制到约化对偶子群对（dual reductive pairs of subgroups）上，先讨论一般情形，再讨论较特殊的情形。

第三部分以数论背景材料开篇，主要讨论 Weil 群（Weil groups）及其 $L$ 函数。随后转向与自守表示相联系的 $L$ 函数、构造它们的各种方法、它们已经证明或仍属猜想的性质，以及有关的局部和整体问题。其余文章大多讨论 $\GL_2$ 的基变换问题（base change problem），以及它在证明某些非阿贝尔 Artin 级数（nonabelian Artin series）的全纯性（holomorphy）方面的应用。

第四部分讨论自守表示与算术代数几何（arithmetical algebraic geometry）之间的联系。在函数域（function fields）情形，它介绍 Drinfeld 关于 $\GL_2$ 的工作；该工作构造出一族 $\ell$ 进表示（$\ell$-adic representations），其 $L$ 级数等于一个给定的自守 $L$ 函数。在数域（number fields）情形，主要讨论 Shimura 簇（Shimura varieties）的问题，包括典范模型（canonical models）、模素理想约化上的点，以及 Hasse--Weil zeta 函数。

本讲习班侧重表示论，因此主要以自守表示的语言来展开讨论。自守表示与自守形式在这里是同一理论的两种表述；整个理论都直接承袭了经典自守形式理论的发展。为了兼顾主题的完整性与日程安排，我们略去了自守形式理论中若干目前尚难纳入这一框架的课题。例如，讲习班在筹划时便希望与 1976 年在波恩举行的单变量模函数会议（Conference on Modular Functions of One Variable）尽量少有重叠。关于其中一些主题及更经典的观点，读者可参阅那次会议的论文集（Springer Lecture Notes 601、627），以及此前几次会议的论文集（Springer Lecture Notes 320、350、476）。此外，约化群、它们的无限维表示或模簇（moduli varieties）等本身颇有意义的主题，在这里也主要依照讲习班核心主题的需要加以讨论。

本论文集分两册出版：第一册包含第一、第二部分，第二册包含第三、第四部分。
\begin{flushright}A. Borel\par W. Casselman\end{flushright}
